% =====================================================================
\chapter{Generación e hiperparámetros}
\label{tema:5}
% =====================================================================

Generar texto es aplicar la regla de la cadena del módulo~1 con el modelo como calculadora de probabilidades condicionadas: se predice la distribución del token siguiente, se muestrea, se añade el token al contexto y se repite. Este tema formaliza el bucle de generación, estudia las tres transformaciones de la distribución que exponen las interfaces de los modelos actuales (temperatura, top-$k$ y top-$p$) y demuestra que la entropía de la distribución crece con la temperatura. Termina con el efecto de los hiperparámetros sobre el entrenamiento del mini-GPT. Se sigue el §6 de la tarea~1 de CS336~\cite{cs336}, el §9.5.5 de~\cite{d2l}, los artículos de Fan, Lewis y Dauphin~\cite{fan} y de Holtzman \emph{et al.}~\cite{holtzman}, y el §11.4 de~\cite{gbc}.

\section{Generación autorregresiva}

\begin{definicion}[Bucle de generación]\label{def:generacion}
Sea un modelo con contexto máximo $T$, una secuencia inicial (\emph{prompt}) $w_{1:m}$, un número máximo de tokens nuevos $G$ y una transformación $\Phi$ que convierte unos logits en una distribución. Para $g=1,\dots,G$:
\begin{enumerate}[label=\arabic*.]
  \item recortar el contexto a sus últimos $T$ tokens, $c=w_{\max(1,\,m+g-T):\,m+g-1}$;
  \item evaluar el modelo sobre $c$ y tomar los logits $z$ de la última posición;
  \item muestrear $w_{m+g}\sim\Phi(z)$ y añadirlo a la secuencia;
  \item detenerse si $w_{m+g}$ es el token delimitador.
\end{enumerate}
\end{definicion}

Por el \cref{cor:una-pasada}, los logits de la última posición son los de la distribución $P_\theta(\cdot\mid c)$. Con $\Phi=\softmax$ se muestrea del propio modelo~\cite[§6, ecuaciones~(21)--(22)]{cs336}; si en cada paso se elige el token más probable, el procedimiento es la \emph{decodificación voraz}~\cite[§9.5.5]{d2l}. Cada paso exige una pasada completa del modelo sobre el contexto; el módulo~19 estudia cómo evitar recalcular lo ya calculado con la caché de claves y valores.

\section{Temperatura}

\begin{definicion}[Softmax con temperatura]\label{def:temperatura}
Para $\tau>0$, $\softmax(z,\tau)_i=e^{z_i/\tau}\big/\sum_j e^{z_j/\tau}$. En este módulo la temperatura se escribe $\tau$ para no confundirla con la longitud de contexto $T$.
\end{definicion}

\begin{proposicion}[Límites]\label{prop:limites}
Sea $z\in\R^{K}$ con un único máximo en el índice $i^{\star}$. Cuando $\tau\to 0^{+}$, $\softmax(z,\tau)\to e_{i^{\star}}$; cuando $\tau\to\infty$, $\softmax(z,\tau)\to(1/K,\dots,1/K)$.
\end{proposicion}

\begin{proof}
Por la \cref{prop:traslacion}, $\softmax(z,\tau)_i=e^{(z_i-z_{i^{\star}})/\tau}\big/\sum_j e^{(z_j-z_{i^{\star}})/\tau}$. Si $\tau\to 0^{+}$, cada exponente con $j\ne i^{\star}$ tiende a $-\infty$ y el de $i^{\star}$ es nulo, luego el denominador tiende a~1 y solo sobrevive el numerador de $i^{\star}$. Si $\tau\to\infty$, todos los exponentes tienden a~0, los numeradores a~1 y el denominador a~$K$.
\end{proof}

\begin{teorema}[La entropía crece con la temperatura]\label{teo:entropia-tau}
Sea $z\in\R^{K}$ no constante y, para $b>0$, $p^{(b)}=\softmax(bz)$, de modo que $b=1/\tau$. Entonces
\[
\frac{d}{db}\,H\bigl(p^{(b)}\bigr)=-b\,\Var_{p^{(b)}}(z)<0,
\]
donde $\Var_{p}(z)=\sum_i p_iz_i^{2}-\bigl(\sum_i p_iz_i\bigr)^{2}$. En consecuencia, $H\bigl(\softmax(z,\tau)\bigr)$ es estrictamente creciente en $\tau$.
\end{teorema}

\begin{proof}
Sea $Z(b)=\sum_j e^{bz_j}$. Como $\ln p^{(b)}_i=bz_i-\ln Z(b)$,
\[
H\bigl(p^{(b)}\bigr)=-\sum_i p^{(b)}_i\bigl(bz_i-\ln Z(b)\bigr)=-b\,\mu(b)+\ln Z(b),\qquad \mu(b)=\sum_i p^{(b)}_iz_i.
\]
Se tiene $\frac{d}{db}\ln Z(b)=\sum_j z_je^{bz_j}/Z(b)=\mu(b)$. Además, $\frac{d}{db}p^{(b)}_i=p^{(b)}_i\bigl(z_i-\mu(b)\bigr)$, luego $\mu'(b)=\sum_i p^{(b)}_iz_i\bigl(z_i-\mu(b)\bigr)=\Var_{p^{(b)}}(z)$. Derivando,
$\frac{d}{db}H=-\mu(b)-b\,\mu'(b)+\mu(b)=-b\Var_{p^{(b)}}(z)$. Como todas las $p^{(b)}_i$ son positivas y $z$ no es constante, la varianza es positiva. Finalmente, $b=1/\tau$ es estrictamente decreciente en $\tau$.
\end{proof}

\begin{ejemplo}\label{ej:temperatura}
Con $z=(2;\allowbreak\,1;\allowbreak\,0{,}5;\allowbreak\,-1)$:
\begin{center}\small
\begin{tabular}{@{}cccccc@{}}
\toprule
$\tau$ & $p_1$ & $p_2$ & $p_3$ & $p_4$ & $H$ (nats) \\
\midrule
$0{,}5$ & $0{,}842$ & $0{,}114$ & $0{,}042$ & $0{,}002$ & $0{,}538$ \\
$1$ & $0{,}609$ & $0{,}224$ & $0{,}136$ & $0{,}030$ & $1{,}014$ \\
$2$ & $0{,}434$ & $0{,}263$ & $0{,}205$ & $0{,}097$ & $1{,}265$ \\
\bottomrule
\end{tabular}
\end{center}
La entropía crece con $\tau$ hacia $\ln 4=1{,}386$, el valor de la distribución uniforme, como afirman la \cref{prop:limites} y el \cref{teo:entropia-tau}.
\end{ejemplo}

La temperatura es la misma que la de la distribución de Boltzmann de la física estadística (módulo~1). Holtzman \emph{et al.} observan que bajar la temperatura mejora la calidad del texto generado a costa de reducir su diversidad~\cite[§3.3]{holtzman}.

\section{Muestreo top-\texorpdfstring{$k$}{k} y top-\texorpdfstring{$p$}{p}}

\begin{definicion}[Top-$k$]\label{def:topk}
Sea $1\le k\le K$ y sea $I_k(z)$ el conjunto de los índices de los $k$ mayores logits, con los empates resueltos por el índice menor. El \emph{muestreo top-$k$} usa la distribución $\softmax\bigl(\topk(z)\bigr)$, donde $\topk(z)_i=z_i$ si $i\in I_k(z)$ y $-\infty$ en otro caso.
\end{definicion}

Fan, Lewis y Dauphin introdujeron este procedimiento con $k=10$ para generar relatos y observaron que producía textos menos repetitivos que la búsqueda en haz~\cite[§5.4]{fan}. Con $k=1$ el muestreo coincide con la decodificación voraz; con $k=K$ no cambia nada. Con $z=(2;\allowbreak\,1;\allowbreak\,0{,}5;\allowbreak\,-1)$, $k=2$ y $\tau=1$, la distribución pasa a ser $(0{,}731;\allowbreak\,0{,}269;\allowbreak\,0;\allowbreak\,0)$: desaparece la cola de tokens improbables, que con 46 caracteres puede acumular una masa apreciable.

\begin{definicion}[Muestreo de núcleo o top-$p$]\label{def:topp}
Sea $q$ una distribución sobre $\V$ y $\rho\in(0,1]$. Sea $V(\rho)$ el menor conjunto de índices tal que $\sum_{j\in V(\rho)}q_j\ge\rho$. El \emph{muestreo de núcleo} usa la distribución
\[
\tilde q_i=\begin{cases} q_i\big/\sum_{j\in V(\rho)}q_j, & i\in V(\rho),\\ 0, & \text{en otro caso}.\end{cases}
\]
\end{definicion}

Es la definición de Holtzman \emph{et al.}~\cite[§3.1, ecuaciones~(2) y~(3)]{holtzman}, reproducida por CS336~\cite[§6, ecuación~(24)]{cs336}, donde el umbral se denota $p$. La distribución $q$ suele ser ya la softmax con temperatura.

\begin{proposicion}[Cálculo del núcleo]\label{prop:nucleo}
Sean $q_{(1)}\ge q_{(2)}\ge\dots\ge q_{(K)}$ las probabilidades ordenadas y $m=\min\{j:\sum_{i\le j}q_{(i)}\ge\rho\}$. Todo conjunto de masa al menos $\rho$ tiene al menos $m$ elementos, y el conjunto $S_m$ de los $m$ índices más probables alcanza esa masa. Por tanto, los conjuntos mínimos tienen exactamente $m$ elementos y $S_m$ es, entre ellos, uno de masa máxima.
\end{proposicion}

\begin{proof}
Un conjunto de $j$ índices tiene masa a lo sumo $\sum_{i\le j}q_{(i)}$, la suma de las $j$ mayores probabilidades. Si $j<m$, esa suma es menor que $\rho$ por la definición de $m$, luego ningún conjunto de menos de $m$ elementos alcanza la masa $\rho$. El conjunto $S_m$ la alcanza, también por la definición de $m$, y su masa es la máxima posible entre los conjuntos de $m$ elementos por la primera observación con $j=m$.
\end{proof}

\begin{observacion}\label{obs:nucleo}
El conjunto mínimo no es único en general, aunque no haya empates: con $q=(0{,}5;\allowbreak\,0{,}3;\allowbreak\,0{,}2)$ y $\rho=0{,}6$, tanto $\{1,2\}$ como $\{1,3\}$ tienen dos elementos y masa al menos $0{,}6$. Por convenio, $V(\rho)$ designa el conjunto $S_m$, que es el que se obtiene al ordenar la distribución y acumular probabilidades hasta alcanzar $\rho$, como prescribe CS336~\cite[§6]{cs336}.
\end{observacion}

La ventaja de top-$p$ sobre top-$k$ es que el número de candidatos se adapta a la forma de la distribución. Con una distribución plana, un $k$ pequeño elimina continuaciones razonables; con una distribución muy concentrada, un $k$ grande deja entrar candidatos inapropiados, cuya probabilidad crece al renormalizar~\cite[§3.2]{holtzman}. Las interfaces de programación de modelos comerciales como Claude y Gemini exponen públicamente los parámetros de temperatura, top-$k$ y top-$p$.

\paragraph{Muestras del mini-GPT.} Con $\tau=0{,}8$ y $k=10$, casi todas las líneas generadas son correctas:
\begin{center}\small
\begin{tabular}{@{}l@{}}
\texttt{objeto. una lámpara azul de cristal. sirve para esperar sentado.}\\
\texttt{objeto. una caja azul de tela. sirve para leer de noche.}\\
\texttt{receta. hornea verde paperado.}
\end{tabular}
\end{center}
La tercera muestra el límite del modelo: mezcla un color propio de los objetos en una receta. Con $\tau=3$ la salida se degrada a símbolos sueltos, porque la distribución se acerca a la uniforme y los dígitos y signos aparecen en cualquier lugar.

\section{Efecto de los hiperparámetros}

El \cref{tab:hiper} resume entrenamientos de 800 pasos sobre el mismo corpus y con la misma semilla, en los que solo cambia lo indicado.

\begin{table}[h]
\centering\small
\caption{Efecto de los hiperparámetros en el mini-GPT (resultados del cuaderno de prácticas). Base: $d=64$, $L=2$, $n_h=4$, $T=32$, $\eta_{\max}=3\cdot 10^{-3}$. Pérdidas de validación en nats por carácter; tiempo relativo al modelo base.}\label{tab:hiper}
\begin{tabularx}{\textwidth}{@{}lrccc>{\raggedright\arraybackslash}X@{}}
\toprule
Cambio & $N$ & Paso 100 & Final & Tiempo & Lectura \\
\midrule
Base & 107\,694 & $2{,}046$ & $0{,}324$ & 1 & referencia \\
$\eta_{\max}=3\cdot 10^{-4}$ & 107\,694 & $2{,}959$ & $1{,}840$ & 1 & demasiado lenta \\
$\eta_{\max}=3\cdot 10^{-2}$ & 107\,694 & $1{,}304$ & $0{,}297$ & 1 & mejor: tolera tasas altas \\
ídem, sin calentamiento & 107\,694 & $1{,}232$ & $0{,}306$ & 1 & algo peor \\
$\eta_{\max}=3\cdot 10^{-1}$ & 107\,694 & $2{,}863$ & $2{,}244$ & 1 & inestable \\
Tasa constante & 107\,694 & $2{,}046$ & $0{,}337$ & 1 & falta el afinado final \\
$d=32$ & 29\,294 & $2{,}308$ & $0{,}649$ & $0{,}5$ & poca capacidad \\
$L=1$ & 57\,902 & $2{,}128$ & $0{,}392$ & $0{,}4$ & un bloque no basta \\
$n_h=1$ & 107\,694 & $2{,}062$ & $0{,}336$ & $0{,}7$ & algo peor \\
$T=8$ & 106\,158 & $2{,}068$ & $0{,}694$ & $0{,}3$ & no ve la etiqueta \\
$d=128$, $L=3$ & 609\,838 & $1{,}441$ & $0{,}297$ & $2{,}4$ & mejor, más caro \\
\bottomrule
\end{tabularx}
\end{table}

La tasa de aprendizaje es el factor dominante: entre $3\cdot 10^{-4}$ y $3\cdot 10^{-2}$ la pérdida final pasa de $1{,}840$ a $0{,}297$, y con $3\cdot 10^{-1}$ el entrenamiento se rompe. Goodfellow, Bengio y Courville la señalan como el hiperparámetro más importante si solo se puede ajustar uno~\cite[§11.4.1]{gbc}, y CS336 pide localizar empíricamente la tasa a partir de la cual el entrenamiento diverge y relacionarla con la mejor~\cite[§7.2, problema \emph{learning\_rate}]{cs336}. El contexto es el segundo factor: con $T=8$ la pérdida se duplica, porque el modelo pierde la información que más reduce la incertidumbre, el tipo de línea (\cref{eje:5-contexto}). Multiplicar por $5{,}7$ los parámetros da la misma mejora que subir la tasa a $3\cdot 10^{-2}$ sin coste adicional. El calentamiento y el coseno aportan poco en este modelo (unas centésimas de nat cada uno), porque es pequeño y el recorte ya lo protege; en los modelos grandes son imprescindibles.

\begin{observacion}[Reglas prácticas]\label{obs:reglas}
De lo anterior se desprenden cuatro reglas: (i) ajustar primero la tasa de aprendizaje, barriendo potencias de 10 o de $\sqrt{10}$, y quedarse con la mayor que sea estable; (ii) antes de escalar, comprobar que el modelo sobreajusta un único lote hasta pérdida casi nula; (iii) elegir $T$ de modo que cubra la dependencia que se quiere capturar, aquí una línea completa de unos 60 a 80 caracteres; (iv) más parámetros solo ayudan si hay datos y pasos suficientes para entrenarlos, como precisan las leyes de escalado del módulo~16.
\end{observacion}

\begin{lecturas}
\item Stanford CS336~\cite{cs336}, tarea~1, §6 «Generating text» (ecuaciones 21--24) y §7.2 «TinyStories».
\item Holtzman \emph{et al.}~\cite{holtzman}, §3 «Language Model Decoding» (§3.1--3.3).
\item Fan, Lewis y Dauphin~\cite{fan}, §5.4 «Generation».
\item Zhang \emph{et al.}~\cite{d2l}, §9.5.5 «Decoding».
\item Goodfellow, Bengio y Courville~\cite{gbc}, §11.4 «Selecting Hyperparameters», en particular §11.4.1.
\item Karpathy~\cite{nanogpt}, archivos \texttt{sample.py} y \texttt{config/train\_shakespeare\_char.py}.
\end{lecturas}

\section*{Ejercicios del tema 5}
\addcontentsline{toc}{section}{Ejercicios del tema 5}

\begin{ejercicio}{1}{eje:5-conceptos}
Describa la distribución de muestreo en cada caso: (a) $k=1$; (b) $k=K$; (c) $\rho=1$; (d) $\tau\to 0^{+}$; (e) $\tau=3$ con 46 caracteres. ¿Cuáles de ellos son deterministas?
\end{ejercicio}

\begin{ejercicio}{2}{eje:5-topk}
Con logits $z=(3,1,1,0)$, aplique top-$k$ con $k=2$ y $\tau=1$. El segundo y el tercer logit están empatados: obtenga la distribución con el convenio de la \cref{def:topk} y con el convenio alternativo que conserva todo logit mayor o igual que el $k$-ésimo mayor.
\end{ejercicio}

\begin{ejercicio}{2}{eje:5-temp}
Con $z=(3,1,0)$, calcule $\softmax(z,\tau)$ y su entropía para $\tau=0{,}5$, $1$ y $2$, y compruebe que la entropía crece con $\tau$. Calcule $\Var_{p}(z)$ para $\tau=1$ y la derivada $dH/db$ en $b=1$ del \cref{teo:entropia-tau}.
\end{ejercicio}

\begin{ejercicio}{2}{eje:5-topp}
La distribución del token siguiente es $q=(0{,}5;\allowbreak\,0{,}2;\allowbreak\,0{,}15;\allowbreak\,0{,}1;\allowbreak\,0{,}05)$. (a) Aplique top-$p$ con $\rho=0{,}8$ y con $\rho=0{,}9$. (b) Aplique antes la temperatura $\tau=0{,}5$ a los logits $\ln q_i$ y después top-$p$ con $\rho=0{,}8$.
\fuente{adaptado de Stanford CS336~\cite{cs336}, tarea~1, §6, problema \emph{decoding} y ecuaciones~(23)--(24) (datos numéricos añadidos).}
\end{ejercicio}

\begin{ejercicio}{3}{eje:5-conmuta}
(a) Demuestre que, para todo $\tau>0$, aplicar top-$k$ y después la temperatura da la misma distribución que aplicar la temperatura y después top-$k$. (b) Demuestre con un contraejemplo que la afirmación análoga para top-$p$ es falsa. (c) ¿Qué orden usa la \cref{def:topp} cuando $q$ es la softmax con temperatura?
\end{ejercicio}

\begin{ejercicio}{3}{eje:5-contexto}
Con $T=8$, el mini-GPT termina con pérdida de validación $0{,}694$ frente a $0{,}324$ con $T=32$ (\cref{tab:hiper}). Use la estructura del corpus del \cref{ej:corpus} para explicar qué información le falta, qué partes de la línea siguen siendo predecibles y cuáles se vuelven inciertas. Cuantifique el exceso de pérdida por línea de 58 caracteres y compárelo con la entropía de la etiqueta.
\end{ejercicio}
